%%%%%%%%%%%%%%%%%%%%%%%%%%%%%%%%%%%%%%%%%%%%%%%%%%%%%%%%%%%%%%%%%%%%%%%%%%%%%%%%
% Obxectivo: Lista de termos empregados no documento,                          %
%            xunto cos seus respectivos significados.                          %
%%%%%%%%%%%%%%%%%%%%%%%%%%%%%%%%%%%%%%%%%%%%%%%%%%%%%%%%%%%%%%%%%%%%%%%%%%%%%%%%

\newglossaryentry{bytecode}{
  name=bytecode,
  description={Código independente da máquina que xeran compiladores de determinadas linguaxes (Java, Erlang,\dots) e que é executado polo correspondente intérprete.}
}

\newglossaryentry{homebrew}{
  name=homebrew,
  description={Software desenvolvido de xeito amateur pola comunidade para unha plataforma, á marxe das canles oficiais do seu fabricante.}
}

\newglossaryentry{modding}{
  name=modding,
  description={Práctica de modificar o hardware orixinal dunha consola, xeralmente para mellorar as súas prestacións ou adaptalo aos estándares actuais, como substituír a pantalla ou engadir unha batería recargable.}
}